\section{Estimación de Dirección de Arribo (DOA)}\label{sec:doa} 

La \textbf{estimación de dirección de arribo (Direction of Arrival, DOA)} busca determinar 
desde qué ángulo llega una onda a un arreglo de sensores.  
Un ejemplo ilustrativo es la ecolocación de los murciélagos: 
pequeñas diferencias de tiempo y fase entre ambos oídos permiten estimar 
la dirección del sonido (véase Fig.~\ref{fig:doa:biologico}).  
En un arreglo de micrófonos ocurre lo mismo: 
cada sensor recibe la misma onda con un retardo relativo, 
y ese desfase contiene la información angular (Fig.~\ref{fig:doa:arregloSensores}) \cite{YangBook}.

\insertarfigura{fig/doaArregloSensores.png}{Arreglo de sensores lineal para estimación de DOA.}{fig:doa:arregloSensores}{1}
\insertarfigura{fig/DoaBiologico.png}{Ejemplo biológico de estimación de dirección de arribo en murciélagos.}{fig:doa:biologico}{1}

\subsection{Modelo de onda plana unidimensional}

Consideremos primero una onda plana armónica propagándose a lo largo de un eje \(x\). 
La señal recibida en la posición \(x\) es:
\begin{equation}
	s(x,t) = A \, e^{j(\beta x - \omega t + \phi_0)},
\end{equation}
donde \(A\) es la amplitud, \(\beta=2\pi/\lambda\) el número de onda, 
\(\omega=2\pi f\) la frecuencia angular y \(\phi_0\) la fase inicial.  

Esta expresión refleja dos aspectos:  
(i) el término temporal \(e^{-j\omega t}\) describe la oscilación sinusoidal,  
(ii) el término espacial \(e^{j\beta x}\) indica cómo cambia la fase al desplazarse en \(x\).  
Cuando la onda avanza una longitud de onda \(\lambda\), 
el exponente completa un ciclo de \(2\pi\), ilustrando su periodicidad (Fig.~\ref{fig:doa:ondaPlana}).  

\insertarfigura{fig/doaOndaPlana.png}{Onda plana en un medio unidimensional.}{fig:doa:ondaPlana}{1}

\subsection{Extensión al caso tridimensional}

En tres dimensiones, la onda plana en \(\mathbf{x} = [x_1,x_2,x_3]^T\) se escribe:
\begin{equation}
	s(\mathbf{x},t) = s(\mathbf{0},t)\, 
	e^{j \beta ( x_1 \cos\varphi \sin\theta + x_2 \sin\varphi \sin\theta + x_3 \cos\theta )},
\end{equation}
donde \(\varphi\) es el \textbf{azimut} y \(\theta\) la \textbf{elevación}.  
Aquí \(s(\mathbf{0},t)\) es la señal de la fuente en su origen.  
En DOA lo importante no es el contenido exacto de \(s(\mathbf{0},t)\), 
sino cómo la propagación introduce diferencias de fase entre sensores \cite{YangBook}.

\subsection{Arreglos linealmente uniformes}

Un caso clásico es el \textbf{Uniform Linear Array (ULA)}: 
una fila de $M$ sensores equiespaciados (Fig.~\ref{fig:doa:arregloSensores}).  
Si la onda incide frontalmente, todos los micrófonos reciben la señal en fase.  
En cambio, con un ángulo oblicuo aparecen retardos progresivos: 
cuanto más inclinado, mayor la diferencia de tiempos entre sensores.  

En términos físicos, cada micrófono mide la misma señal con una \textbf{fase distinta}.  
Ese patrón de desfases codifica la dirección de llegada.  
Matemáticamente, la señal en el micrófono $m$-ésimo es:
\begin{equation}
	s_m(t) = s(0,t)\, a_m(\varphi,\theta),
\end{equation}
donde $s(0,t)$ es la señal de referencia y $a_m(\varphi,\theta)$ un factor de fase:
\begin{equation}
	a_m(\varphi,\theta) = e^{j \beta m d \cos\varphi \sin\theta}.
\end{equation}

\subsection{Steering vector}

Agrupando todos los factores de fase en un vector columna se obtiene el \textbf{steering vector}:
\begin{equation}
	\mathbf{a}(\varphi,\theta) =
	\begin{bmatrix}
		1 \\
		e^{j \beta d \cos\varphi \sin\theta} \\
		e^{j 2 \beta d \cos\varphi \sin\theta} \\
		\vdots \\
		e^{j (M-1)\beta d \cos\varphi \sin\theta}
	\end{bmatrix}.
\end{equation}

Este vector es la \textbf{firma espacial} del arreglo: 
si la onda proviene de \((\varphi,\theta)\), las fases se alinean, 
y los algoritmos de DOA (DAS, Capon, MUSIC, ESPRIT) lo utilizan como base 
para estimar la dirección de la fuente \cite{YangBook,Gupta2015}.  




\subsection{Delay-and-Sum Beamformer (DAS)}\label{sec:doa:das}

El método más clásico y sencillo para estimar la dirección de arribo es el 
\textbf{beamformer de retraso y suma} (DAS, por sus siglas en inglés).  
La idea es intuitiva: si se supone un ángulo de llegada, se compensan los retardos 
con los que la onda alcanza a cada sensor, y luego se suman las señales.  

\begin{itemize}
	\item Cuando el ángulo supuesto coincide con el real, las señales quedan alineadas en fase 
	y se refuerzan al sumarse (suma constructiva).  
	\item Si el ángulo no corresponde, las señales llegan desfasadas 
	y se atenúan entre sí (suma destructiva).  
\end{itemize}

De esta forma, basta recorrer un conjunto de ángulos candidatos y calcular la potencia 
de la salida del beamformer para obtener el \textbf{espectro espacial DAS}: 
los máximos de este espectro indican las direcciones estimadas de llegada de la señal 
\cite{YangBook}.

\subsubsection*{Características principales}
\begin{itemize}
	\item Implementación extremadamente simple y de bajo costo computacional.  
	\item No requiere información adicional sobre el entorno ni sobre las señales.  
	\item Resolución angular limitada: los lóbulos principales son anchos, 
	lo que dificulta separar fuentes cercanas.  
	\item Rechazo a interferencias poco eficiente.  
	\item Mejora con mayor número de sensores y mayor apertura del arreglo.  
\end{itemize}

\subsubsection*{Formulación matemática}

Sea $\mathbf{u}(t) \in \mathbb{C}^M$ el vector de observaciones de un arreglo de $M$ sensores, 
y $R_{uu}=\mathbb{E}\{\mathbf{u}(t)\mathbf{u}^H(t)\}$ su matriz de correlación.  
El \emph{steering vector} para un ángulo candidato $\varphi$ es:

\begin{equation}
	\mathbf{a}(\varphi) =
	\begin{bmatrix}
		1 \\
		e^{j 2 \pi d_{\text{norm}} \cos \varphi} \\
		e^{j 2 \pi d_{\text{norm}} \cdot 2 \cos \varphi} \\
		\vdots \\
		e^{j 2 \pi d_{\text{norm}} (M-1)\cos \varphi}
	\end{bmatrix},
\end{equation}

donde $d_{\text{norm}} = d/\lambda$ es el espaciamiento normalizado entre elementos.  

La salida del beamformer apuntando hacia $\varphi$ es:
\begin{equation}
	y(t) = \mathbf{a}^H(\varphi)\, \mathbf{u}(t),
\end{equation}

y su potencia promedio:
\begin{equation}
	P_{DAS}(\varphi) = \mathbf{a}^H(\varphi)\, R_{uu}\, \mathbf{a}(\varphi).
\end{equation}

Al graficar $P_{DAS}(\varphi)$ en función del ángulo, los picos corresponden 
a las direcciones estimadas de arribo.  
En la Fig.~\ref{fig:doa:dasMatlab} se muestra un ejemplo de este espectro 
simulado en MATLAB para un arreglo lineal.

\insertarfigura{fig/DAS.png}{Espectro espacial obtenido con el beamformer Delay-and-Sum (DAS). 
	Se simularon tres fuentes incidentes en $30^\circ$, $50^\circ$ y $70^\circ$, 
	con una relación señal-ruido (SNR) de 5~dB, usando un arreglo de 10 sensores 
	uniformemente espaciados con $d_{norm}=0,5$. 
	Los picos del espectro corresponden a las direcciones estimadas de arribo.}{fig:doa:dasMatlab}{1}
 \subsection{Capon Beamformer (MVDR)}\label{sec:doa:capon}

El beamformer de \textbf{Capon}, también conocido como 
\textbf{MVDR (Minimum Variance Distortionless Response)}, 
se propuso como una mejora directa del método Delay-and-Sum.  
Su idea central es apuntar hacia una dirección deseada sin distorsionar la señal proveniente de allí, 
al mismo tiempo que se minimiza la energía recibida desde todas las demás direcciones.  

De manera intuitiva, si el DAS se asemeja a una linterna con un haz ancho, 
Capon actúa como un láser mucho más fino: concentra la ganancia únicamente en el ángulo de interés 
y suprime la contribución de señales interferentes \cite{YangBook}.

\begin{itemize}
	\item \textbf{Ganancia unitaria en la dirección deseada}: se garantiza que la señal en el ángulo $\varphi$ no se distorsione.  
	\item \textbf{Mínima potencia desde otras direcciones}: los pesos del arreglo se ajustan para cancelar la energía procedente del resto del espacio.  
\end{itemize}

\subsubsection*{Características}
\begin{itemize}
	\item Mayor resolución que DAS: lóbulos principales más estrechos y mejor discriminación angular.  
	\item Rechazo eficaz de interferencias fuera de la dirección de interés.  
	\item Requiere una estimación precisa de la matriz de correlación $R_{uu}$; con pocos datos puede degradarse.  
	\item Más costoso computacionalmente, ya que exige invertir matrices.  
	\item Sensible a errores de calibración en el arreglo.  
\end{itemize}

\subsubsection*{Formulación matemática}

El problema de diseño de Capon se plantea como una optimización:
\begin{equation}
	\min_{\mathbf{w}} \; \mathbf{w}^H R_{uu}\, \mathbf{w}
	\quad \text{sujeto a} \quad \mathbf{w}^H \mathbf{a}(\varphi) = 1,
\end{equation}
donde $R_{uu}$ es la matriz de correlación de las señales, 
$\mathbf{a}(\varphi)$ es el \emph{steering vector}, 
y $\mathbf{w}$ los pesos del beamformer.

La solución, obtenida con multiplicadores de Lagrange, es:
\begin{equation}
	\mathbf{w}_c(\varphi) = \frac{R_{uu}^{-1}\, \mathbf{a}(\varphi)}
	{\mathbf{a}^H(\varphi)\, R_{uu}^{-1}\, \mathbf{a}(\varphi)}.
\end{equation}

El espectro espacial asociado al beamformer de Capon es:
\begin{equation}
	P_{\text{Capon}}(\varphi) = 
	\frac{1}{\mathbf{a}^H(\varphi)\, R_{uu}^{-1}\, \mathbf{a}(\varphi)}.
\end{equation}

Los máximos de $P_{\text{Capon}}(\varphi)$ señalan las direcciones estimadas de arribo, 
con lóbulos mucho más estrechos que los de DAS, 
tal como se aprecia en la Fig.~\ref{fig:doa:caponMatlab}.

\insertarfigura{fig/CAPON.png}{Espectro espacial obtenido con el beamformer de Capon (MVDR). 
	Se simularon tres fuentes incidentes en $30^\circ$, $50^\circ$ y $70^\circ$, 
	con una relación señal-ruido (SNR) de 5~dB, usando un arreglo de 10 sensores 
	con espaciamiento normalizado $d_{norm}=0,5$. 
	En comparación con DAS, los lóbulos principales son más estrechos, 
	lo que permite una mejor discriminación angular entre fuentes cercanas.}{fig:doa:caponMatlab}{1}
 \subsection{MUSIC (Multiple Signal Classification)}\label{sec:doa:music}

El algoritmo \textbf{MUSIC} es uno de los métodos más influyentes para la estimación de dirección de arribo.  
Su principal fortaleza es la capacidad de alcanzar \textbf{alta resolución}, separando fuentes muy cercanas angularmente incluso con un número moderado de sensores \cite{Gupta2015}.  

La idea se basa en descomponer la matriz de correlación $R_{uu}$ en dos subespacios ortogonales:
\begin{itemize}
	\item El \textbf{espacio de señal}, generado por los \emph{steering vectors} de las fuentes reales.
	\item El \textbf{espacio de ruido}, ortogonal al espacio de señal.
\end{itemize}

Un steering vector correspondiente a una dirección real de llegada es ortogonal al subespacio de ruido.  
Por lo tanto, al proyectar $\mathbf{a}(\varphi)$ sobre ese subespacio, la proyección se anula en los ángulos correctos.

\subsubsection*{Características}
\begin{itemize}
	\item Alta resolución, capaz de distinguir fuentes muy próximas.  
	\item Requiere conocer o estimar el número de fuentes presentes.  
	\item Sensible a correlaciones entre señales y a errores de calibración del arreglo.  
\end{itemize}

\subsubsection*{Formulación matemática}

La descomposición en valores propios de $R_{uu}$ permite separar:
\begin{equation}
	R_{uu} = E_s \Lambda_s E_s^H + E_n \Lambda_n E_n^H,
\end{equation}
donde $E_s$ contiene los autovectores asociados a las fuentes (espacio de señal) y $E_n$ los del ruido.

El espectro espacial de MUSIC se define como:
\begin{equation}
	P_{\text{MUSIC}}(\varphi) = 
	\frac{1}{\mathbf{a}^H(\varphi)\, E_n E_n^H\, \mathbf{a}(\varphi)}.
\end{equation}

Los picos de $P_{\text{MUSIC}}(\varphi)$ corresponden a las direcciones estimadas de llegada.  

\insertarfigura{fig/MUSIC.png}{Espectro espacial obtenido con el algoritmo MUSIC. 
	Se simularon tres fuentes incidentes en $30^\circ$, $50^\circ$ y $70^\circ$, 
	con SNR de 5~dB y un arreglo ULA de 10 sensores con $d_{norm}=0,5$. 
	Los picos estrechos permiten distinguir claramente las fuentes.}{fig:doa:musicMatlab}{1}

\subsection{Root-MUSIC}\label{sec:doa:rootmusic}

El método \textbf{Root-MUSIC} es una variante diseñada para evitar el barrido angular.  
En lugar de evaluar $P_{\text{MUSIC}}(\varphi)$ punto a punto, 
transforma el problema en encontrar raíces de un polinomio complejo asociado al subespacio de ruido.  

\subsubsection*{Ventajas}
\begin{itemize}
	\item Reduce la complejidad computacional (no requiere barrido angular).  
	\item Ofrece estimaciones más precisas en comparación con MUSIC clásico.   
\end{itemize}

\subsubsection*{Limitación principal}
\begin{itemize}
	\item Está restringido a arreglos lineales uniformes (ULA). 
\end{itemize}

\subsubsection*{Formulación matemática}

Se define el polinomio:
\begin{equation}
	Q[z] = \sum_{k=-(M-1)}^{M-1} q_k z^{-k}, 
	\qquad q_k = \sum_{m-l=k} C(m,l),
\end{equation}
con $z = e^{j 2\pi \tfrac{d}{\lambda} \cos\varphi}$ y $C = E_n E_n^H$.  

Las raíces $z_i$ cercanas al círculo unitario determinan los ángulos estimados:
\begin{equation}
	\hat{\varphi}_i = \cos^{-1}\!\left(\frac{\angle z_i}{2\pi d}\right).
\end{equation}

\subsection{Spatial Smoothing MUSIC}\label{sec:doa:ssmusic}

El \textbf{spatial smoothing} es una extensión de MUSIC para manejar señales \emph{coherentes} (por ejemplo, en presencia de multitrayectoria).  
Cuando las señales son coherentes, la matriz $R_{uu}$ pierde rango y MUSIC clásico falla.  
El smoothing restaura el rango promediando matrices de correlación de subarreglos \cite{Gupta2015}.  

\subsubsection*{Características}
\begin{itemize}
	\item Permite detectar señales coherentes que MUSIC clásico no puede resolver.  
	\item A costa de reducir la apertura efectiva del arreglo (menor resolución).  
	\item \emph{Forward smoothing}: detecta hasta $M/2$ señales coherentes.  
	\item \emph{Forward/Backward smoothing}: hasta $2M/3$ señales coherentes.  
\end{itemize}

\subsubsection*{Formulación matemática}

Sea $K=M-p+1$ el número de subarreglos forward posibles.  
La correlación promedio es:
\begin{equation}
	R^f = \frac{1}{K}\sum_{k=0}^{K-1} R^f_k.
\end{equation}

Si además se consideran subarreglos backward, se obtiene:
\begin{equation}
	R = \frac{R^f + R^b}{2}.
\end{equation}

\insertarfigura{fig/MUSIC-SS.png}{Ejemplo de MUSIC con spatial smoothing (SS-MUSIC). 
	Se consideran tres fuentes coherentes en $30^\circ$, $50^\circ$ y $70^\circ$, 
	con SNR de 5~dB y 10 sensores ($d_{norm}=0,5$). 
	El smoothing restaura el rango de la matriz de correlación y permite separar correctamente las fuentes.}{fig:doa:ssmusicMatlab}{1}

\insertarfigurawide{fig/comparativoDOA.png}{Comparación de métodos de estimación de DOA: 
	DAS muestra lóbulos anchos y baja resolución; 
	Capon mejora la nitidez pero aún con interferencias; 
	MUSIC separa claramente las fuentes en condiciones adversas; 
	y ESPRIT obtiene los ángulos sin necesidad de barrido.}{fig:doa:comparacion}{1}
 \subsection{ESPRIT (Estimation of Signal Parameters via Rotational Invariance Techniques)}\label{sec:doa:esprit}

El algoritmo \textbf{ESPRIT} es una alternativa a MUSIC que evita el barrido angular.  
En lugar de calcular un espectro y buscar picos, se apoya en la 
\textbf{invariancia rotacional} entre subarreglos de un ULA (Uniform Linear Array).  
Esto permite estimar directamente las direcciones de arribo resolviendo un problema algebraico \cite{Gupta2015}.

\subsubsection*{Idea general}
\begin{itemize}
	\item Se divide el arreglo de $M$ sensores en dos subarreglos idénticos, 
	desplazados por una distancia fija $\Delta x$.  
	\item La señal en el segundo subarreglo es equivalente a la del primero, 
	multiplicada por un factor de fase dependiente del ángulo de arribo.  
	\item Esta relación se modela como una \textbf{invariancia rotacional}, 
	lo que permite estimar los ángulos a partir de los autovalores de una matriz reducida.  
\end{itemize}

Mientras que MUSIC “escanea” ángulo por ángulo para detectar picos,  
ESPRIT utiliza propiedades algebraicas del arreglo y evita la búsqueda explícita.

\subsubsection*{Características}
\begin{itemize}
	\item No requiere barrido angular, reduciendo la complejidad frente a MUSIC.  
	\item Menor tiempo de cálculo y menos sensible a discretización de ángulos.  
	\item Exige que el arreglo posea subestructuras desplazadas (ej. ULA).  
	\item Sensible a errores de calibración o desviaciones geométricas.  
\end{itemize}

\subsubsection*{Formulación matemática}

Dividiendo el arreglo en dos subarreglos, las señales medidas se modelan como:
\begin{equation}
	\mathbf{u}_0(t) = A\, \mathbf{s}(t) + \mathbf{n}_0(t), 
	\qquad 
	\mathbf{u}_1(t) = A \Phi \mathbf{s}(t) + \mathbf{n}_1(t),
\end{equation}
donde $A$ es la matriz de steering del primer subarreglo, 
y $\Phi$ una matriz diagonal de factores de fase:
\[
\Phi = \mathrm{diag}\left(e^{j\gamma_1}, e^{j\gamma_2}, \dots, e^{j\gamma_L}\right), 
\qquad \gamma_l = \frac{2\pi}{\lambda}\, \Delta x \cos\varphi_l.
\]

A partir del subespacio de señal $V_s$, se extraen dos bloques $V_0$ y $V_1$ 
correspondientes a cada subarreglo.  
La relación de invariancia se expresa como:
\begin{equation}
	V_1 \approx V_0 \Psi,
\end{equation}
donde $\Psi$ es una matriz de rotación que contiene la información angular.

Los autovalores de $\Psi$ son:
\[
\exp(j \gamma_l) = e^{j \tfrac{2\pi}{\lambda}\, \Delta x \cos\varphi_l},
\]
a partir de los cuales se obtiene directamente la estimación de los ángulos:
\begin{equation}
	\hat{\varphi}_l = \cos^{-1}\!\left( \frac{\lambda}{2\pi \Delta x} \, \angle(e^{j \gamma_l}) \right), 
	\quad l=0,\dots,L-1.
\end{equation}


Finalmente, para contrastar los distintos métodos se simuló un escenario común 
(Fig.~\ref{fig:doa:comparacion}) con tres fuentes en $30^\circ$, $50^\circ$ y $70^\circ$, 
relación señal-ruido de 5~dB, un arreglo ULA de 10 sensores con espaciamiento normalizado $d_{norm}=0,5$, 
y multitrayectoria (trayectoria directa más una retardada).  
En el caso de MUSIC con \emph{spatial smoothing}, se emplearon $M-2$ subarreglos solapados.


\subsection{Relevancia y aplicaciones}

Los métodos de \textbf{estimación de dirección de arribo (DOA)} responden a una necesidad central: 
\textbf{localizar espacialmente una fuente sonora o electromagnética a partir de un arreglo de sensores}.  
En el análisis de audio, esto significa no solo detectar que un mosquito está presente, 
sino también \textbf{saber desde qué dirección llega su zumbido}, lo que habilita aplicaciones de seguimiento 
y control (por ejemplo, orientar un dron hacia el objetivo).

El problema que se aborda es: \textit{¿cómo separar y ubicar múltiples señales simultáneas en un entorno ruidoso?}.  
Los algoritmos clásicos como DAS o Capon proporcionan soluciones basadas en filtrado espacial, 
fáciles de interpretar pero limitadas en resolución.  
Técnicas como MUSIC y ESPRIT aprovechan la descomposición en subespacios de la matriz de covarianza 
para alcanzar una \textbf{resolución angular muy superior}, incluso cuando las fuentes están cercanas.  
Cada enfoque es más adecuado según el contexto: métodos clásicos ofrecen bajo costo computacional 
y robustez, mientras que los subespaciales brindan mayor precisión a costa de mayor complejidad.

Más allá de los mosquitos, los algoritmos de DOA se aplican en un amplio rango de problemas:  
\begin{itemize}
	\item \textbf{Comunicaciones inalámbricas}: beamforming en 5G/6G para dirigir haces adaptativos.  
	\item \textbf{Acústica arquitectónica}: ubicar fuentes de ruido en edificios o ciudades inteligentes.  
	\item \textbf{Sonar y radar}: localizar embarcaciones, submarinos o aeronaves en entornos adversos.  
	\item \textbf{Biomedicina}: identificar focos de actividad en el cerebro a partir de EEG/MEG.  
\end{itemize}

En todos estos casos, la lógica es la misma: aprovechar la información espacial contenida en las diferencias 
de fase y amplitud entre sensores para estimar ángulos de llegada.  
Así, la estimación de DOA se convierte en un paso natural que extiende la simple detección de señales, 
aportando un componente espacial que es indispensable para sistemas de localización y seguimiento modernos.


